\section{Exact endpoint-path parity waves}
\label{sec:direct-path}

The one-edge transient is not the only exactly solvable case. Consider the
path with vertices $0,1,\ldots,n-1$, seeded at endpoint $0$, and suppose the
far endpoint is not reached in the argument below. Thus $d_0=1$ and
$d_i=2$ on every displayed vertex $i\geq1$. Define the normalized residual
and two scalar coefficients
\begin{equation}
 y_i:=\frac{r_i}{\sqrt{d_i}},
 \qquad
 \lambda:=\lambda_\alpha,
 \qquad
 \eta:=\frac{\lambda}{1+\lambda^2}.
 \label{eq:direct-path-normalization}
\end{equation}
The endpoint seed gives $y^{(0)}_0=\alpha$. An optimal-SOR push at $u$ has
the particularly simple form
\begin{equation}
 y_u^+=-\lambda^2y_u,
 \qquad
 y_v^+=y_v+\frac{2\lambda}{d_v}y_u
 \quad(v\sim u).
 \label{eq:direct-path-sor-update}
\end{equation}
Thus an interior target receives $\lambda y_u$, whereas endpoint $0$
receives $2\lambda y_u$. In these coordinates the empirical rank is
\begin{equation}
 p_u(y)=\frac{d_u}{1+d_u}|y_u|
 =\begin{cases}|y_0|/2,&u=0,\\2|y_u|/3,&u\geq1.
 \end{cases}
 \label{eq:direct-path-rank}
\end{equation}

Put $A_m:=\alpha\lambda^m$. For $m\geq0$, define the signed plateau
\begin{equation}
 S_m:\qquad
 y_i=\begin{cases}
  A_m,&0\leq i\leq m\text{ and }i\equiv m\pmod2,\\
  -A_{m+1},&0\leq i\leq m\text{ and }i\not\equiv m\pmod2,\\
  0,&i>m.
 \end{cases}
 \label{eq:direct-path-plateau}
\end{equation}
The first state is exactly $S_0$.

\begin{lemma}[One parity sweep advances one layer]
\label{lem:direct-path-sweep}
Starting from $S_m$, push once every coordinate
\begin{equation*}
 H_m:=\{i:0\leq i\leq m,\ i\equiv m\pmod2\}
\end{equation*}
with optimal SOR. The pushes commute, and the resulting state is $S_{m+1}$.
\end{lemma}

\begin{proof}
The vertices in $H_m$ are nonadjacent. A pushed value $A_m$ becomes
$-\lambda^2A_m=-A_{m+2}$. Every old low-parity interior coordinate receives
two packets $\lambda A_m$, so its value changes as
\begin{equation*}
 -\lambda A_m+2\lambda A_m=A_{m+1}.
\end{equation*}
When endpoint $0$ is low parity it has only one pushed neighbor, but the
endpoint coefficient is $2\lambda$, giving the same identity. Finally,
vertex $m+1$ receives one packet $\lambda A_m=A_{m+1}$. These are precisely
the entries of $S_{m+1}$.
\end{proof}

The sweeps need not be executed level by level. Introduce the event set
\begin{equation}
 \mathcal E:=\{(m,u):m\geq0,\ 0\leq u\leq m,\ u\equiv m\pmod2\},
 \label{eq:direct-path-events}
\end{equation}
where event $(m,u)$ means pushing coordinate $u$ at positive value $A_m$.
The event's dependencies are
\begin{equation}
\begin{aligned}
 \operatorname{pred}(m,m)&=\{(m-1,m-1)\},
 &&m\geq1,\\
 \operatorname{pred}(m,0)&=\{(m-1,1)\},
 &&m\geq2\text{ even},\\
 \operatorname{pred}(m,u)&=\{(m-1,u-1),(m-1,u+1)\},
 &&0<u<m.
\end{aligned}
\label{eq:direct-path-predecessors}
\end{equation}

\begin{lemma}[Corner-flip representation]
\label{lem:direct-path-corner-flip}
After any dependency-respecting finite event prefix, let $K_u$ be the number
of executed events at coordinate $u$ and define its next-event height
\begin{equation}
 m_u:=u+2K_u.
 \label{eq:direct-path-height}
\end{equation}
Then
\begin{equation}
 |m_{u+1}-m_u|=1
 \label{eq:direct-path-one-lipschitz}
\end{equation}
and, for every interior coordinate,
\begin{equation}
 y_u=\begin{cases}
  +A_{m_u},&m_{u-1}=m_u+1=m_{u+1},\\
  -A_{m_u},&m_{u-1}=m_u-1=m_{u+1},\\
  0,&m_{u-1}\neq m_{u+1}.
 \end{cases}
 \label{eq:direct-path-corner-residual}
\end{equation}
At endpoint $0$, the first case means $m_1=m_0+1$ and the second means
$m_1=m_0-1$. Thus positive residuals are exactly local minima of the height
interface, negative residuals are local maxima, and a legal push is the
corner flip $m_u\leftarrow m_u+2$.
\end{lemma}

\begin{proof}
Initially $K_u=0$ and $m_u=u$, so adjacent heights differ by one. The next
event at $u$ is $(m_u,u)$. The predecessor rules
\cref{eq:direct-path-predecessors} say exactly that it is enabled when $m_u$
is a local minimum. Executing it increments $K_u$, hence raises $m_u$ by
two and turns that minimum into a maximum while preserving
\cref{eq:direct-path-one-lipschitz}. The three residual values follow from
the packet accounting used in \cref{lem:direct-path-sweep}: zero, one, or two
executed interior predecessors leave respectively
$-A_{m_u},0,+A_{m_u}$, while the endpoint's single doubled predecessor gives
the two endpoint cases. Induction proves the claim.
\end{proof}

\begin{lemma}[Every negative packet has a larger positive certificate]
\label{lem:direct-path-positive-certificate}
After any dependency-respecting prefix of events, each coordinate's next
event has one of three states:
\begin{equation}
 +A_m\ \text{(enabled)},\qquad
 -A_m\ \text{(waiting)},\qquad
 0\ \text{(one interior predecessor received)}.
 \label{eq:direct-path-three-states}
\end{equation}
If $y_u=-A_m$, then for some
\begin{equation*}
 1\leq j\leq\frac{m-u}{2}
\end{equation*}
the northeast event $(m-j,u+j)$ is enabled, is at an interior coordinate,
and has positive residual $A_{m-j}>A_m$.
\end{lemma}

\begin{proof}
The three states are the three cases of
\cref{eq:direct-path-corner-residual}. If $y_u=-A_m$, then $m_u=m$ is a local
maximum and $m_{u+1}=m-1$. Walk right along the height interface. The height
cannot decrease forever: the event prefix is finite, so $K_v=0$ and $m_v=v$
for all sufficiently large $v$. Let $u+j$ be the first coordinate where the
slope changes from $-1$ to $+1$. It is a local minimum with height $m-j$ and
hence, by \cref{lem:direct-path-corner-flip}, has positive residual
$A_{m-j}$. Also $m_{u+j}=m-j\geq u+j$ by
\cref{eq:direct-path-height}, so $j\leq(m-u)/2$. Finally $j\geq1$ and
$u+j>0$, making the certificate interior. Since $0<\lambda<1$, its amplitude
$A_{m-j}$ is strictly larger than $A_m$.
\end{proof}

\begin{lemma}[Refreshed absolute ranking is an event schedule]
\label{lem:direct-path-ranked-events}
Start from $S_0$ and, after every operation, choose a live coordinate of
maximum empirical rank \cref{eq:direct-path-rank}, using absolute residuals
exactly as in the measured policy. Every selected residual is positive and
equals $A_m$ for a unique event $(m,u)$. Moreover, before the first inactive
event level, the executed events form a dependency-respecting ordering of
\cref{eq:direct-path-events}; the conclusion is independent of tie breaking.
\end{lemma}

\begin{proof}
By \cref{lem:direct-path-positive-certificate}, a negative interior entry
$-A_m$ has score $2A_m/3$ and is accompanied by a positive interior entry
with score $2A_{m-j}/3>2A_m/3$. A negative endpoint entry has the still
smaller score $A_m/2$ and the same type of interior certificate. If the
negative entry is active, its larger certificate is active as well.
Therefore an absolute-rank maximizer can never be negative.

By \cref{eq:direct-path-three-states}, a positive entry is exactly an enabled
event and has its prescribed amplitude. The selected operation is therefore
legal. Induction from $S_0$ proves the dependency-respecting event statement.
There are finitely many active event levels, so the maximizer executes every
event above the gate.
\end{proof}

\begin{theorem}[Exact refreshed two-rung theorem on an endpoint path]
\label{thm:direct-path-two-rung}
Assume $0<\alpha<1$, $1<B<B_{\rm edge}$, and $B\epsppr<1$. Let
\begin{equation}
 L:=\min\{m\geq0:\lambda^m\leq B\epsppr\}
 =\left\lceil\frac{\log(B\epsppr)}{\log\lambda}\right\rceil,
 \label{eq:direct-path-depth}
\end{equation}
and take $n\geq L+4$. The fully refreshed absolute-score-maximizing spreading
phase ends exactly at plateau $S_L$ and performs
\begin{equation}
 N_{\rm spread}
 =\sum_{m=0}^{L-1}\left(\left\lfloor\frac m2\right\rfloor+1\right)
 =\left\lfloor\frac{(L+1)^2}{4}\right\rfloor
 \label{eq:direct-path-push-count}
\end{equation}
coordinate pushes. Its exact charged work is
\begin{equation}
 \Work_{\rm spread}
 =3N_{\rm spread}-\left\lceil\frac L2\right\rceil.
 \label{eq:direct-path-spread-work}
\end{equation}

The fully refreshed terminal exact phase pushes every active coordinate in
the high-parity set $H_L$ at most once and possibly pushes coordinate $L+1$
once. It then certifies the terminal gate. In particular,
\begin{equation}
 \Work_{\rm exact}
 \leq3\left(\left\lfloor\frac L2\right\rfloor+2\right).
 \label{eq:direct-path-exact-work}
\end{equation}
\end{theorem}

\begin{proof}
The spreading gate is $B\tau=B\alpha\epsppr$. By
\cref{lem:direct-path-ranked-events}, precisely the events with $A_m$ above
that gate execute. These are the levels $0,\ldots,L-1$. Completing those
events in any dependency-respecting order gives $S_L$ by
\cref{lem:direct-path-sweep} and linearity. Level $m$ contains
$\lfloor m/2\rfloor+1$ events, proving
\cref{eq:direct-path-push-count}. Every interior push costs three and there
is one cost-two endpoint event at each even level, proving
\cref{eq:direct-path-spread-work}.

It remains to analyze the exact rung. Write $A:=A_L$. If $A\leq\tau$, it is
already finished. Otherwise every coordinate in $H_L$ is active. Under an
exact push, an interior target receives $\eta y_u$ and endpoint $0$ receives
$2\eta y_u$. The high-parity coordinates are nonadjacent and retain value
$A$ until pushed. While an unpushed high-parity interior coordinate remains,
its score $2A/3$ exceeds every low-parity score. If endpoint $0$ is the last
high-parity coordinate, its score $A/2$ still exceeds all remaining scores:
the only adjacent partially canceled value has magnitude
$A\lambda^3/(1+\lambda^2)$, and
\begin{equation*}
 \frac23\frac{\lambda^3}{1+\lambda^2}A\leq\frac A3<\frac A2.
\end{equation*}
Thus refreshed ranking pushes each member of $H_L$ once before any low
coordinate.

After these pushes, every old low-parity coordinate has magnitude
\begin{equation}
 R(\lambda)A,
 \qquad
 R(\lambda):=\frac{\lambda(1-\lambda^2)}{1+\lambda^2}
 \leq R_\star,
 \label{eq:direct-path-terminal-ratio}
\end{equation}
where the endpoint receives one doubled packet and hence has the same value
as an interior coordinate receiving two packets. The new outer coordinate
$L+1$ has value $\eta A$. Since $BR_\star<1$, all old low-parity entries are
below $\tau$. If $\eta A>\tau$, coordinate $L+1$ is the only active entry
and is pushed once. It leaves packets of magnitude $\eta^2A$ at $L$ and
$L+2$. Finally,
\begin{equation*}
 B\eta^2\leq B/4<1
\end{equation*}
because $B<B_{\rm edge}<4$. Hence those packets are also below $\tau$.
There are $\lfloor L/2\rfloor+1$ high-parity vertices, all charges in this
part are at most three, and \cref{eq:direct-path-exact-work} follows.
\end{proof}

The explored prefix has degree volume
$\mathcal V_{\rm exp}=1+2L$ and charge volume $2+3L$. Since
\begin{equation}
 -\log\lambda=2\operatorname{arctanh}(\sqrt\alpha)
 \geq2\sqrt\alpha,
 \label{eq:direct-path-log-lambda}
\end{equation}
the theorem immediately gives
\begin{equation}
 \Work
 =O\!\left(\mathcal V_{\rm exp}
 \left[1+\frac{\log(1/(B\epsppr))}{\sqrt\alpha}\right]\right).
 \label{eq:direct-path-explored-volume-bound}
\end{equation}
The exact count is $\Theta(L^2)=\Theta(\mathcal V_{\rm exp}L)$ for this
trajectory. Thus the bound has the desired accelerated propagation length
$1/\sqrt\alpha$ per exposed volume, with the logarithm needed to attenuate
from the seed scale to the chosen gate. It is an algorithm-specific path
identity, not a graph-uniform oracle lower bound.

\begin{corollary}[Sign-guarded literal batches preserve the path theorem]
\label{cor:direct-path-sign-guarded-batches}
Under the assumptions of \cref{thm:direct-path-two-rung}, use the literal
top-fraction snapshots and arbitrary sequential batch size,
but strengthen each live guard on the endpoint path to require both the
current threshold and $r_u>0$. Keep the same empirical absolute-residual
snapshot ordering. Then every charged spreading operation is a legal event,
every nonempty spreading batch executes at least one event, and the spreading
phase has exactly \cref{eq:direct-path-push-count,eq:direct-path-spread-work}.
If the same sign guard is used in the exact rung, its charged sequence obeys
\cref{eq:direct-path-exact-work} and certifies the terminal gate.
\end{corollary}

\begin{proof}
At the start of any spreading batch, the largest snapshot score is positive
by \cref{lem:direct-path-positive-certificate}; it is the first operation and
therefore passes the live sign guard. Every later charged coordinate also has
positive live residual, so \cref{eq:direct-path-three-states} identifies it
as an enabled event. Skipped negative or zero entries cost no coordinate work.
Thus every batch extends a legal event prefix and makes progress. There are
finitely many events above the gate, proving the spreading claims.

At plateau $S_L$, the ranking comparison in the proof of
\cref{thm:direct-path-two-rung} places an unpushed positive high-parity entry
before every negative low-parity entry. High-parity entries are nonadjacent,
so each batch executes at least one and all charged high-parity operations are
unchanged legal exact pushes. Selected negative entries are skipped. After
the high parity is settled, all old low-parity entries are below the gate;
the only possible active entry is the positive outer coordinate $L+1$, which
is pushed once. The previous terminal bound applies.
\end{proof}

Under the declared coordinate-operation meter, the extra sign test adds no
charged work. It can increase the number of frontier snapshots or service
round trips, which is a separate engineering cost and is not hidden in the
corollary.

For the two values used in the measured study, the sign guard can in fact be
removed.  The essential fact is that the endpoint's smaller charge weight
cannot make it fall behind the next amplitude level.

\begin{theorem}[Literal top-$1/32$ path theorem in the measured regime]
\label{thm:direct-path-literal-measured-regime}
Under the endpoint-path assumptions of
\cref{thm:direct-path-two-rung}, suppose additionally that
\begin{equation}
 \lambda<\frac34,
 \qquad\text{equivalently}\qquad
 \alpha>\frac1{49}.
 \label{eq:direct-path-literal-alpha-window}
\end{equation}
Run the original unguarded absolute-residual policy with literal batch size
\begin{equation*}
 k=\min\{8192,\lceil|\mathcal F_\theta|/32\rceil\}.
\end{equation*}
Then every charged operation in both rungs is positive and legal,
independently of tie breaking among equal-rank entries.  The literal method
has exactly the spreading event multiset and work
\cref{eq:direct-path-push-count,eq:direct-path-spread-work}, and its terminal
rung obeys \cref{eq:direct-path-exact-work}.  In particular, the theorem
covers the measured values $\alpha=0.025$ and $0.04$ without adding a sign
guard.
\end{theorem}

\begin{proof}
First consider the spreading rung.  Induct on batches using the following
invariant: for some event level $m$, every earlier level is complete and an
arbitrary subset $X\subseteq H_m$ has executed.  Let
$U:=|H_m\setminus X|$ and let $C$ be the number of nonzero low-parity
coordinates between the sites of $H_m$.  The unexecuted high sites have
value $+A_m$, the executed high sites have value $-A_{m+2}$, and a low site
has value $-A_{m+1}$, $0$, or $+A_{m+1}$ according as zero, one, or all of
its level-$m$ predecessors have executed.

Encode $X$ as a binary word on the ordered sites of $H_m$.  A maximal run of
$r$ executed sites creates $r-1$ active internal low sites.  Charge the one
remaining site of the run to the following unexecuted high site, or, if the
run reaches the outermost high site, to the new active outer low site.  The
charges are disjoint, so
\begin{equation}
 |X|\leq U+C.
 \label{eq:direct-path-partial-layer-count}
\end{equation}
If height $m+2$ is inactive, the frontier has at most $U+C$ entries.  If it
is active, every lower nonzero stratum is active and
\begin{equation*}
 |\mathcal F_\theta|=U+|X|+C\leq2(U+C).
\end{equation*}
In either case, whenever the frontier is nonempty,
\begin{equation}
 k\leq U+C.
 \label{eq:direct-path-batch-does-not-reach-self-negative}
\end{equation}
Thus a snapshot cannot reach the self-negative height-$m+2$ stratum.

For any height $h$, \cref{eq:direct-path-literal-alpha-window} gives
\begin{equation}
 \frac12A_h>\frac23A_{h+1}.
 \label{eq:direct-path-cross-height-rank}
\end{equation}
Hence every height-$h$ entry, even the endpoint, strictly outranks every
height-$(h+1)$ entry.  If the batch reaches height $m+1$, every remaining
height-$m$ positive occurs earlier in its sequential list.  A selected
positive height-$(m+1)$ entry was already enabled.  A selected negative one
had none of its predecessors at the snapshot; all of those height-$m$
predecessors execute earlier, changing it to zero or to the positive legal
event before its turn.  Therefore every charged operation is legal.  The
batch ends either at another partial level $m$ or at a partial level $m+1$,
closing the induction.  Gate truncation only deletes a suffix of amplitude
levels and the cap can only decrease $k$.

For the exact rung, write $A=A_L$.  After $s=0,1,2$ predecessor pushes, an
old low-parity entry has respectively the values
\begin{equation}
 -\lambda A,
 \qquad
 -\frac{\lambda^3}{1+\lambda^2}A,
 \qquad
 R(\lambda)A>0,
 \label{eq:direct-path-exact-partial-values}
\end{equation}
where the endpoint's one doubled predecessor corresponds directly to the
last case.  An unpushed high entry remains $+A$ and a pushed high entry is
zero.  By \cref{eq:direct-path-cross-height-rank}, even an endpoint high
entry outranks an untouched interior low entry; the one-packet magnitude in
\cref{eq:direct-path-exact-partial-values} is smaller still.  Thus if a
snapshot reaches a negative low entry, all of its remaining high
predecessors occur earlier and make it positive or inactive.  No negative
exact push is charged.

After every high site has been pushed, all old lows have magnitude at most
$R_\star A<\tau$ because $BR_\star<1$.  The only possible live coordinate
is the positive outer site with value $\eta A$.  Its one exact push leaves
packets bounded by $\eta^2A<\tau$, since $B\eta^2<1$.  This is precisely the
terminal count and certificate of \cref{thm:direct-path-two-rung}.
\end{proof}

\begin{proposition}[Unguarded full-frontier batching fails on the path]
\label{prop:direct-path-full-batch-fails}
Use the original live absolute-threshold guard, but select the complete
frontier in every spreading batch.  If the gate is below $A_4$, then on the
endpoint path the third batch executes a negative-residual push.  Hence the
unguarded analogue of
\cref{cor:direct-path-sign-guarded-batches} is false for arbitrary batch
sizes.
\end{proposition}

\begin{proof}
Represent the state by the corner heights from
\cref{lem:direct-path-corner-flip}.  The first batch pushes endpoint $0$ and
changes the initial heights near the endpoint from
\begin{equation*}
 (0,1,2,3,4)\quad\hbox{to}\quad(2,1,2,3,4).
\end{equation*}
The second snapshot contains the positive minimum $1$ at height one and the
negative endpoint $0$ at height two.  Their ranks satisfy
\begin{equation*}
 \frac23A_1>\frac12A_2,
\end{equation*}
so the batch pushes coordinate $1$ first.  This enables coordinate $0$,
which then also executes legally.  The resulting heights are
\begin{equation*}
 (4,3,2,3,4).
\end{equation*}
The third snapshot contains the positive minimum $2$ at height two and the
negative endpoint $0$ at height four, ordered as
\begin{equation*}
 \frac23A_2>\frac12A_4.
\end{equation*}
Pushing coordinate $2$ does not change the nonadjacent endpoint.  Because
$A_4$ remains above the gate, the live absolute-threshold guard accepts
coordinate $0$ next and applies optimal SOR to its negative residual.  This
is not a legal corner-flip event.
\end{proof}

Thus the measured fraction $1/32$ is not an inessential implementation
constant in this proof program.  A sufficiently small snapshot must prevent
a negative maximum from outrunning the chain of positive minima that enables
it.  The next proposition isolates the exact closure property needed from
that fraction.

\begin{proposition}[Parent-closed batch coupling]
\label{prop:direct-path-parent-closed-batch}
Consider a top-$k$ snapshot taken from a legal path event-frontier state.
Suppose every selected negative waiting coordinate has at least one of its
direct positive predecessor events selected earlier in the sequential batch.
Then every selected coordinate is either skipped by the live guard or
executes a legal positive event. Consequently, a sequence of parent-closed
batches has exactly the refreshed event multiset of
\cref{lem:direct-path-ranked-events}.
\end{proposition}

\begin{proof}
A negative interior entry waiting for event $(m,u)$ equals $-A_m$. Each
selected direct predecessor executes earlier at amplitude $A_{m-1}$ and
contributes $A_m$. One such predecessor changes the entry to zero, so it is
skipped unless the other predecessor also executes first; two change it to
$+A_m$, enabling precisely event $(m,u)$. At endpoint $0$, its single
predecessor contributes $2A_m$ and changes $-A_m$ directly to $+A_m$. Direct
interior predecessors have strictly larger rank than the waiting entry; the
hypothesis explicitly requires the relevant predecessor to occur earlier in
the declared order. Positive enabled events are nonadjacent and do not stale
one another. Induction down the selected list proves the first statement, and
induction over batches proves the second.
\end{proof}

The parent-closure hypothesis cannot be proved from the geometry of an
arbitrary legal corner interface alone.  The next exact construction is
reachable by positive legal events but its top-$1/32$ snapshot is not
parent-closed.  Thus the history invariant in
\cref{thm:direct-path-literal-measured-regime} is essential.

\begin{proposition}[Static top-$1/32$ parent closure is false]
\label{prop:direct-path-static-parent-closure-fails}
Take $\lambda=1/2$, equivalently $\alpha=1/9$, and use spreading gate
$\theta=A_{39}$.  There is a dependency-respecting legal event prefix whose
frontier has $34$ coordinates and whose top-two snapshot first selects the
positive coordinate $2$ and then charges a still-negative endpoint.
\end{proposition}

\begin{proof}
Define the event counts
\begin{equation*}
 K_0=18,
 \qquad K_1=17,
 \qquad K_i=16\quad(2\leq i\leq6),
\end{equation*}
and, for $0\leq s\leq15$, define
\begin{equation*}
 K_{7+2s}=K_{8+2s}=15-s,
 \qquad K_i=0\quad(i\geq39).
\end{equation*}
This sequence is finitely supported and nonincreasing, with every adjacent
drop equal to zero or one.  Its event ideal
\begin{equation*}
 \mathcal I(K):=\{(u+2t,u):0\leq t<K_u\}
\end{equation*}
is predecessor-closed: a left predecessor uses
$K_{u-1}\geq K_u$, while a right predecessor uses
$K_{u+1}\geq K_u-1$.  Executing the ideal in increasing event height
therefore reaches the interface using only positive legal events.

The heights $m_u=u+2K_u$ begin
\begin{equation*}
 (36,35,34,35,36,37,38),
\end{equation*}
followed by alternating corners of heights $37,38$ and then the increasing
tail $39,40,\ldots$.  The $34$ active corners are exactly
\begin{equation*}
 y_0=-A_{36},
 \qquad y_2=+A_{34},
\end{equation*}
together with sixteen negative entries $-A_{38}$ at
$u=6,8,\ldots,36$ and sixteen positive entries $+A_{37}$ at
$u=7,9,\ldots,37$.  Hence the literal batch size is two, and its ranks obey
\begin{equation*}
 \frac23A_{34}>\frac12A_{36}>\frac23A_{37}
 =\frac13A_{36}.
\end{equation*}
The snapshot uniquely orders coordinate $2$ before endpoint $0$.  The first
push changes only coordinates $1$ and $3$, so the endpoint remains
$-A_{36}$, is still strictly above $\theta$, and is charged negatively.
\end{proof}

This proposition does not say that the literal trajectory from $S_0$
reaches this particular interface.  It says that no static theorem over all
legal reachable interfaces can prove the desired coupling.  For
$\alpha>1/49$, the partial-layer history invariant supplies the missing
restriction.  For smaller $\alpha$, the endpoint can lag multiple interior
amplitude levels and an all-parameter literal path theorem remains open.

\paragraph{Computed diagnostic (not theorem).}
Across $28$ path instances with four gate scales and
\begin{equation*}
 \alpha\in\{0.2,0.1,0.04,0.01,0.0025,0.000625,10^{-4}\},
\end{equation*}
the literal $1/32$ batches selected $4{,}603$ negative snapshot entries.
Every one satisfied the parent-closure condition and no negative entry was
charged.  This is consistent with the proved partial-layer invariant in the
measured regime and with, but does not prove, its all-$\alpha$ extension.

\begin{remark}[What remains literal]
The measured implementation selects a top fraction from one snapshot and
then applies live guards. Direct simulations of the path recurrence produce
the same event multiset and terminal plateau: selected negative waiting
coordinates are either canceled and skipped or enabled before their turn.
The refreshed theorem removes snapshot staleness, while
\cref{thm:direct-path-literal-measured-regime} handles it exactly for
$\alpha>1/49$.  The static counterexample
\cref{prop:direct-path-static-parent-closure-fails} shows why the same proof
cannot be extended merely by considering arbitrary legal corner interfaces.
For $\alpha\leq1/49$, simulations still produce the refreshed event multiset,
but a stronger trajectory-history invariant or an actual-run counterexample
is needed.
\end{remark}
